% Generated draft: requires mathematical and source review.
\documentclass[11pt]{article}
\usepackage[margin=25mm]{geometry}
\usepackage{fontspec}
\setmainfont{Latin Modern Roman}
\usepackage{amsmath,amssymb,mathtools,unicode-math}
\setmathfont{Latin Modern Math}
\usepackage{enumitem,needspace,xurl,hyperref}
\setlength{\parindent}{0pt}
\setlength{\parskip}{5pt}
\setlength{\emergencystretch}{3em}
\newcommand{\N}{\mathbb{N}}
\newcommand{\Z}{\mathbb{Z}}
\newcommand{\Q}{\mathbb{Q}}
\newcommand{\R}{\mathbb{R}}
\newcommand{\C}{\mathbb{C}}
\newcommand{\F}{\mathbb{F}}
\newcommand{\A}{\mathbb{A}}
\newcommand{\PP}{\mathbb{P}}
\newcommand{\E}{\mathbb{E}}
\newcommand{\eps}{\varepsilon}
\newcommand{\dd}{\,\mathrm{d}}
\begin{document}
\section*{UnsolvedMath ID: 7600021}
\textbf{Source problem code:} catalogue entry
\subsection{Definitions and mathematical statement}
For a structure \(A\) in a specified effective first-order language, \(\operatorname{Th}(A)\) is decidable if an algorithm determines, for every sentence in that language, whether \(A\) satisfies it. The conventional field/ring language is \(\{0,1,+,-,\cdot\}\); naming a parameter such as \(t\), or additional constants or a valuation, changes the formal question.

The first source question concerns ``the field of Laurent series over \(\mathbb Z_p\)''. For a coefficient ring \(A\),
\[
 A((t))=\left\{\sum_{n\ge N}a_nt^n:N\in\mathbb Z,\ a_n\in A\right\}.
\]
If \(\mathbb Z_p\) means \(\mathbb Z/p\mathbb Z=\mathbb F_p\) for a prime \(p\), this is the field \(\mathbb F_p((t))\). If it means the ring of \(p\)-adic integers, \(\mathbb Z_p((t))\) is a ring rather than a field. Replacing its coefficients by \(\mathbb Q_p\) would give a different field and is not done here. The notation and quantification over \(p\) must therefore be resolved before a single decidability problem is identified.

The second question says ``field of polynomials over \(\mathbb C\)''. The polynomial object \(\mathbb C[t]\) is not a field; its fraction field is \(\mathbb C(t)\). The supplied summary explicitly interprets this part as asking whether
\[
 \operatorname{Th}(\mathbb C(t))\text{ is decidable}.
\]
That interpretation is retained conditionally, not equated with decidability of \(\mathbb C[t]\) or the Laurent-series field \(\mathbb C((t))\).

The summary quotes Scanlon describing the \(\mathbb C(t)\) question as long-standing open, mentions an unfinished Pheidas undecidability strategy, and attributes decidability of \(\mathbb C((t))\) to Ax--Kochen in connection with R. Robinson's question. It leaves the first part unverified and ends mid-sentence. No exact language or complete bibliographic details are supplied. These source reports do not resolve the notation or establish independently verified current status; primary-source review remains outstanding.
\subsection{Short English statement}
Ask for decidability of the intended Laurent-series theory and, on the summary's reading, of \(\mathbb C(t)\). The original wording confuses coefficient-ring notation and polynomial rings with fields, and the language and prime quantifiers must be specified.
\subsection{Sources}
\begin{itemize}
\item Local registry: \url{data/unsolvedmath/problems.json}, permanent UnsolvedMath ID \(7600021\), source code catalogue entry; supplied statement, background, and research summary, updated 2026-08-06.
\item Wikipedia contributors, \emph{List of unsolved problems in mathematics}, Model theory and formal languages bullet 21, pinned revision 1366636767: \url{https://en.wikipedia.org/w/index.php?title=List_of_unsolved_problems_in_mathematics&oldid=1366636767#Model_theory_and_formal_languages}. MediaWiki-source extraction accessed 2026-07-29; primary-source status audit remains required.
\item Scanlon, ``Decidability of some complicated structures definable in \(\mathbb C(t)\)'', quoted by the supplied summary; no full publication details are supplied.
\item Pheidas; Ax--Kochen; and R. Robinson are named by the summary without exact references. Its first-part assessment is visibly truncated.
\end{itemize}
\end{document}
